\chapter{Conclusion}

In this thesis, we introduced Diffold, a diffusion-based generative framework for RNA three-dimensional structure prediction that integrates (i) the geometric backbone and evolutionary representations of RhoFold+, (ii) structural priors from the large-scale RNA language model RNA-FM, and (iii) an Elucidated Diffusion Model (EDM) formulation for continuous-noise coordinate generation. By directly modeling all-atom coordinates through an iterative denoising process conditioned on residue- and atom-level features, Diffold aims to improve not only global fold recovery but also local geometric consistency and physical plausibility.

Comprehensive evaluation on the combined CASP16 and RNA-benchmark test set demonstrates that Diffold consistently surpasses the strong baseline RhoFold+ across all major accuracy and geometry metrics. Specifically, Diffold reduces RMSD from $3.85 \pm 1.96$ $\text{\AA}$ to $2.29 \pm 1.53$ $\text{\AA}$ (−40.5\%) and improves TM-score from $0.309 \pm 0.185$ to $0.585 \pm 0.328$ (+89.3\%), indicating substantially better global fold reconstruction. Consistently, GDT-TS increases from $0.313 \pm 0.262$ to $0.537 \pm 0.365$ (+71.5\%), further supporting improved backbone-level agreement with native structures. For local structural fidelity, Diffold achieves a higher lDDT ($0.845 \pm 0.247$ vs. $0.806 \pm 0.155$, +4.9\%), suggesting improved local geometric consistency on many targets. Importantly, Diffold also yields a markedly lower clash score ($0.0041 \pm 0.0030$ vs. $0.0079 \pm 0.0026$, −48.1\%), reflecting fewer severe steric conflicts and overall better physical plausibility of the generated all-atom conformations. When placed in a broader context using TM-score distributions across representative baselines, Diffold exhibits the highest median TM-score and remains competitive in spread, highlighting the advantage of diffusion-based generation under this evaluation setting.

Beyond aggregate accuracy, Diffold shows improved data independence under a training–test similarity diagnostic based on structural TM-score. Using the maximum TM-score between each test target and all training structures as a proxy for the closest training neighbor, Diffold exhibits only weak and non-significant dependence between test performance and training-set similarity $r = −0.209, p = (1.2\times10^{-1})$, whereas RhoFold+ shows a strong positive correlation $r = 0.752, p = 3.5\times10^{-11}$. This contrast suggests that Diffold’s diffusion generator is less reliant on template-like proximity to structurally similar training examples and can better generalize to structurally more distant targets under this protocol.

Finally, we demonstrated that a lightweight force-field-based post-processing step can effectively correct local covalent-geometry pathologies that may arise when generating coordinates directly in Cartesian space. In the 8UYG case (135 nt), the diffusion output exhibited pronounced backbone breaks due to abnormally elongated $\text{O3}'–\text{P}$ phosphodiester bonds; the number of $\text{O3}'–\text{P}$ distances exceeding 2.0 $\text{\AA}$ decreased from 75 to 1 after AMBER-based relaxation, while preserving the overall fold. This observation supports AMBER relax as a practical post-processing step to improve chemical completeness and physical plausibility when explicit bond constraints are not enforced during generation.

\subsubsection{Discussion}

The results in Chapter 4 suggest that Diffold improves RNA 3D structure prediction accuracy over the RhoFold+ baseline on the combined CASP16 and RNA-benchmark evaluation. Across global and local quality metrics, Diffold achieves substantially lower RMSD and higher TM-score/GDT-TS, while also reducing steric clashes and slightly improving lDDT. This pattern is consistent with the design goal of Diffold: instead of producing a single deterministic structure through direct regression, the diffusion-based generator explicitly models a distribution over plausible conformations and iteratively refines noisy coordinates toward a physically consistent configuration under sequence- and evolution-derived conditions. In practice, this iterative denoising paradigm appears to be particularly effective at correcting large-scale global fold errors that would otherwise dominate RMSD and depress TM-score, while still maintaining local geometry as reflected by lDDT and clash score. Importantly, the gain is not isolated to a single metric; improvements co-occur across multiple evaluation criteria, supporting the claim that Diffold’s advantage reflects a genuine improvement in structural fidelity rather than metric-specific artifacts.

A key implication of these observations is that conditioning a generative diffusion process on strong upstream representations can be more sample-efficient than learning the entire folding mapping end-to-end from scratch. In Diffold, the RhoFold+ backbone provides residue-level single and pair representations derived from evolutionary signals and geometric constraints, and the diffusion module focuses on learning a robust inverse-noising operator in all-atom coordinate space. This division of labor may help explain why the improvements are pronounced even under a training setup that freezes the backbone: the backbone supplies a rich, structured prior over contacts and long-range dependencies, while the diffusion module contributes the ability to explore and refine conformations in a way that is less brittle than one-shot coordinate prediction. From a methodological perspective, this is encouraging because it indicates that diffusion can serve as a general-purpose “structure generator head” that upgrades an existing folding trunk, rather than requiring a full redesign of the upstream architecture.

The post-processing results further clarify the distinction between global structural accuracy and chemical validity. Because Diffold generates full-atom coordinates directly, it does not explicitly enforce ideal covalent geometry for every bond and angle during generation. The observed elongation of the $\text{O3}'–\text{P}$ phosphodiester bond in a subset of predictions is therefore a plausible failure mode: even if the overall fold is correct, localized bond distortions can appear as the model prioritizes global geometric consistency under the learned distribution. Applying AMBER relax as a force-field-based energy minimization step effectively repairs these local backbone discontinuities, dramatically reducing the number of abnormal $\text{O3}'–\text{P}$ distances in the representative case while leaving global evaluation metrics essentially unchanged. This supports an important practical conclusion: the relax step is best interpreted as a chemical “sanity correction” that improves physical realism and molecular integrity, rather than a mechanism that inflates benchmark scores. In other words, Diffold’s accuracy improvements primarily come from the learned generative process, whereas relaxation mainly resolves local covalent-geometry violations that are not fully captured by the diffusion objective.

The length-dependent analysis provides additional insight into where Diffold gains come from and why its behavior may differ from conventional expectations. In many prediction pipelines, performance tends to decrease with increasing sequence length because longer RNAs introduce more long-range constraints, more alternative folds, and larger conformational spaces. However, Diffold shows a positive correlation between sequence length and TM-score in the tested set. There are several plausible explanations. First, the evaluation set may be length-imbalanced: short RNAs can be structurally diverse and may include flexible hairpins, small riboswitch fragments, or motifs with multiple near-degenerate conformations, which are intrinsically difficult to resolve from sequence alone. Longer RNAs in CASP-style benchmarks often include more stable and information-rich folds with stronger evolutionary constraints, and therefore may provide clearer signals to the upstream trunk and the diffusion conditioner. Second, Diffold’s training data are capped at 256 nt, so performance on sequences substantially longer than this cap can be viewed as an out-of-distribution generalization test. The fact that Diffold remains accurate for a 452-nt target while the baseline collapses suggests that the diffusion module may be leveraging backbone-derived global constraints in a way that scales more gracefully with length, at least for certain fold classes. Finally, metric properties may also contribute: TM-score is length-normalized and is relatively forgiving to localized deviations when the global topology is correct; thus, if Diffold better preserves overall topology on long RNAs, TM-score can increase even when some local regions remain imperfect.

The data-independence analysis strengthens the interpretation that Diffold’s performance is not simply explained by memorization or training-set leakage. Using the maximum TM-score between each test structure and the training set as a similarity proxy, the observed weak correlation for Diffold indicates that high test performance does not require the presence of near-identical training exemplars. In contrast, the stronger positive dependence observed for RhoFold+ suggests that the baseline benefits more directly from proximity to training structures, which is a known concern when training data are harvested from PDB without careful family-level partitioning. While this analysis cannot fully rule out subtle biases, it provides evidence that the diffusion generator behaves more like a generalizable conditional model rather than a retrieval-like system. In practical terms, this makes Diffold more attractive for challenging targets where homologous structural templates are scarce or where the fold family is underrepresented in PDB-derived training sets.

The case studies illustrate both the strengths and current limitations of Diffold. The 8uyg example demonstrates that relaxation is an effective post-processing step for repairing backbone continuity when $\text{O3}'–\text{P}$ bond lengths become abnormal, highlighting that local chemical integrity remains a distinct objective from global fold accuracy in diffusion-based coordinate generation. The R1283v2/9j3r example is more striking: Diffold nearly reproduces the native structure with an exceptionally high TM-score on a 452-nt RNA, whereas RhoFold+ fails catastrophically. This suggests that Diffold’s iterative denoising can rescue cases where deterministic prediction may become trapped in poor local minima or propagate early backbone errors throughout the structure. At the same time, the 7wia failure case cautions that diffusion does not automatically solve all hard instances, even when training-set similarity is high. Several factors could plausibly contribute: the target may contain idiosyncratic motifs, pseudoknots, or ligand-/ion-dependent stabilizations that are not captured by the current conditioning features; the MSA quality for that sequence might be poor despite structural similarity to training examples; the structure may include regions with alternative conformations or unresolved experimental density that confuse supervision; or the diffusion sampling could have higher variance for this target, making a single sampled structure unrepresentative. This type of failure also motivates a practical improvement specific to diffusion models: generating multiple samples per target and selecting the best candidate using confidence predictions (e.g., pLDDT/PAE) could reduce the risk of unlucky samples, which is less relevant for deterministic baselines.

\subsubsection{Limitations}
Despite these encouraging results, several limitations should be emphasized.  
First, the training data follow the RhoFold+ pipeline and are derived from PDB representative sets; while this ensures a fair comparison with prior methods, it also implies that the training distribution may not fully reflect the diversity of biologically relevant RNAs, and performance may be biased toward well-studied structural classes.

Second, although redundancy between training and test structures is carefully controlled through clustering and sequence- and structure-level proximity analysis, this procedure does not completely exclude all possible forms of information leakage. In particular, the RNA language model RNA-FM employed in this work was pretrained on large-scale RNA sequence corpora whose overlap with the evaluation targets is not explicitly constrained. As a result, sequence-level information related to test RNAs may have been implicitly encoded in the pretrained representations. Therefore, the observed performance gains should be interpreted as strong empirical improvements rather than a strict demonstration of generalization under a fully leakage-free setting.

Third, the reliance on remotely queried multiple sequence alignments introduces additional biological context at inference time. While this setting reflects practical RNA structure prediction scenarios, it differs from a strictly single-sequence prediction regime. Consequently, the reported accuracy should be understood as conditional on the availability and quality of homologous sequence information, and variability in external database contents may also affect reproducibility over time.

From a computational perspective, diffusion models are inherently more expensive than one-shot predictors in both training and inference. Moreover, the current workflow requires an additional AMBER relaxation step to enforce local covalent-geometry correctness, indicating that the generative model alone does not yet fully guarantee chemically valid backbones.

Finally, the present evaluation focuses on RNA monomers, whereas many biologically critical RNAs function in complexes with proteins, ions, or small-molecule ligands. Extending the framework to these settings will likely require more explicit conditioning on chemical environment and interaction partners.

Overall, these limitations delineate clear directions for future work, including stricter control of pretraining data overlap, improved evaluation under single-sequence settings, tighter integration of chemical constraints into diffusion generation, and extensions to heterogeneous RNA–protein or RNA–ligand complexes.
